% ============================================================
% CHAPTER 18: SEMIDEFINITE PROGRAMMING (matches Vazirani Ch. 18)
% ============================================================
\setcounter{chapter}{17}
\chapter{Semidefinite Programming}
\label{ch:sdp_maxcut}

\section{Intuition: High-Dimensional Embeddings and Hyperplanes}

\begin{intuitionbox}
\textbf{Peeling and Layering Intuition:} The Max-Cut problem is NP-hard and cannot be approximated within $16/17 \approx 0.941$ unless P=NP. However, we can obtain a $0.878$ approximation using Semidefinite Programming (SDP). We model the vertices as unit vectors on a sphere in $\mathbb{R}^n$. Instead of assigning each vertex to 0 or 1, we embed them such that connected vertices are pushed as far apart as possible (minimizing their dot products). To round these vectors, we choose a random hyperplane. Vertices on one side of the hyperplane form one partition, and the rest form the other.
\end{intuitionbox}

\section{Problem Definition}
Given an undirected graph $G = (V, E)$ with non-negative edge weights $w_{ij} \ge 0$, the \emph{Maximum Cut} (Max-Cut) problem is to find a partition of the vertices $V$ into two sets $S$ and $V \setminus S$ such that the total weight of edges crossing the partition is maximized.

The problem is NP-hard. We seek a randomized $0.878$-approximation using the Goemans-Williamson algorithm~\cite{goemans-williamson94}.

\section{Algorithm and Python Implementation}
We implement a lightweight gradient projection optimizer to find the vector embeddings on the unit sphere and apply randomized hyperplane rounding to partition the vertices.

\begin{lstlisting}[style=python, caption=Goemans-Williamson Max-Cut Algorithm]
import math
import random


def compute_cut_weight(assignment, edges, weights):
    """Total weight of edges crossing the cut."""
    return sum(w for (u, v), w in zip(edges, weights)
               if assignment[u] != assignment[v])


def optimize_max_cut_vectors(n, edges, weights, dim=8, lr=0.05, epochs=200):
    """Finds unit vector embeddings in R^dim for Max-Cut."""
    v = []
    for _ in range(n):
        vec = [random.gauss(0.0, 1.0) for _ in range(dim)]
        mag = math.sqrt(sum(x*x for x in vec))
        v.append([x/mag for x in vec])
        
    for _ in range(epochs):
        new_v = []
        for i in range(n):
            grad = [0.0] * dim
            for (u, val_v), weight in zip(edges, weights):
                if u == i:
                    for k in range(dim): grad[k] += weight * v[val_v][k]
                elif val_v == i:
                    for k in range(dim): grad[k] += weight * v[u][k]
            updated = [v[i][k] - lr * grad[k] for k in range(dim)]
            mag = math.sqrt(sum(x*x for x in updated))
            new_v.append(v[i] if mag < 1e-9 else [x/mag for x in updated])
        v = new_v
    return v

def goemans_williamson_max_cut(n, edges, weights, vectors, trials=500):
    """Applies randomized hyperplane rounding to partition vertices."""
    dim = len(vectors[0])
    best_weight, best_assignment = -1.0, []
    for _ in range(trials):
        r = [random.gauss(0.0, 1.0) for _ in range(dim)]
        mag = math.sqrt(sum(x*x for x in r))
        r = [x/mag for x in r]
        
        assignment = [1 if sum(vectors[i][k] * r[k] for k in range(dim)) >= 0 else 0 for i in range(n)]
        weight = compute_cut_weight(assignment, edges, weights)
        if weight > best_weight:
            best_weight, best_assignment = weight, assignment
    return best_assignment, best_weight
\end{lstlisting}

\section{Analysis: The 0.878 Performance Bound}
We formulate the Max-Cut problem as a quadratic integer program: maximize $\frac{1}{2} \sum_{(i,j) \in E} w_{ij} (1 - y_i y_j)$ s.t. $y_i \in \{-1, 1\}$.
By relaxing the scalar variables $y_i \in \{-1, 1\}$ to unit vectors $v_i \in \mathbb{R}^n$, we obtain the Semidefinite Programming relaxation:
\[
\text{Maximize} \quad \frac{1}{2} \sum_{(i,j) \in E} w_{ij} (1 - v_i \cdot v_j) \quad \text{s.t.} \quad v_i \cdot v_i = 1 \; (\forall i \in V)
\]
Under Goemans and Williamson~\cite{goemans-williamson94} randomized rounding, we choose a random unit vector $r$ and partition based on $\text{sign}(v_i \cdot r)$. The probability that edge $(i, j)$ is cut (i.e. $v_i, v_j$ are separated by the random hyperplane) is exactly $\theta_{ij}/\pi$, where $\theta_{ij} = \arccos(v_i \cdot v_j)$ is the angle between $v_i$ and $v_j$.
By linearity of expectation, the expected weight of the cut is:
\[
E[\text{Weight}] = \sum_{(i,j) \in E} w_{ij} \frac{\theta_{ij}}{\pi}
\]
Comparing this term-by-term with the SDP relaxation's objective, we bound the ratio:
\[
\alpha = \min_{\theta \in [0, \pi]} \frac{\theta/\pi}{(1 - \cos \theta)/2} \approx 0.87856
\]
Thus, the expected cut weight is at least $0.878 \cdot \text{OPT}_{SDP} \ge 0.878 \cdot \text{OPT}$.

\section{Applications}
\begin{itemize}
    \item \textbf{Statistical Physics (Ising Spin Glasses):} Finding the minimum energy configuration (ground state) of spin glass magnetic systems.
    \item \textbf{VLSI Layout and Cross-Talk Minimization:} Splitting circuit gates into two layers to minimize inter-layer connection crossovers.
    \item \textbf{Image Segmentation:} Clustering pixels into foreground and background by maximizing cut weight along contrast boundaries.
\end{itemize}

\subsection*{Real Example: VLSI Circuit Cross-Talk Minimization}
\textbf{Scenario:} A chip designer (e.g. Intel) layouts $n = 5,000$ transistors on a chip. Certain pairs of wires have high electrical interference (cross-talk) if routed on the same layer. The designer wants to partition transistors into two separate routing layers to minimize same-layer cross-talk.
\textbf{Model:} Max-Cut. Transistors are vertices, cross-talk interferences are edge weights. Layer division is the cut. The Goemans-Williamson algorithm partitions the circuit layers.
\textbf{Why SDP matters:} Deterministic algorithms easily get stuck in local minima on highly interconnected networks. The Goemans-Williamson SDP algorithm embeds the netlist geometrically, placing highly interfering nodes on opposite sides of a sphere, and partitions them globally with a guaranteed $87.8\%$ efficiency.

\section{Example: The 5-Cycle Graph}
Consider the 5-cycle graph $C_5$ with unit edge weights. The optimal Max-Cut value is 4 (achieved by partitioning into sets of size 2 and 3).
The optimal 2D embedding of $C_5$ places the vectors on a circle with angle $\theta = 4\pi/5 = 144^\circ$ between adjacent nodes.
The SDP objective value is:
\[
\text{OPT}_{SDP} = 5 \cdot \frac{1 - \cos(144^\circ)}{2} \approx 5 \cdot 0.9045 = 4.5225
\]
The expected cut weight under random hyperplane rounding is:
\[
E[\text{Weight}] = 5 \cdot \frac{144^\circ}{180^\circ} = 4.0000
\]
So on this instance the ratio of the expected cut to the SDP value is $4.0 / 4.5225 \approx 0.8844$. This is \emph{above} the worst-case constant $\alpha \approx 0.87856$: the $5$-cycle gives a specific attainable ratio but does \emph{not} attain the exact minimum. The constant $\alpha$ arises as the infimum of $\theta/\pi \div (1-\cos\theta)/2$; it is approached in the limit by a family of instances whose angles cluster near the minimum of that function, not by $C_5$ itself.

\section{Exercises}
\begin{exercise}
Prove that the probability that two unit vectors $v_i, v_j$ are separated by a random hyperplane $r$ is exactly $\theta_{ij}/\pi$, where $\theta_{ij} = \arccos(v_i \cdot v_j)$.
\end{exercise}
\begin{exercise}
Show that for a triangle graph $K_3$ with unit weights, the optimal Max-Cut is 2, while the SDP relaxation value is 2.25. What is the ratio?
\end{exercise}

